\section{Discussion}
\label{sec:discussion}

\paragraph{The signal matters more than the parameters.}~The most
consequential finding of this study is that KEDA's advantage is
rooted in \emph{what} it measures rather than in \emph{how} the
underlying HPA control block is configured. Both scalers were given
an identical $60$\,s stabilisation window, an identical
\texttt{Percent 50\%/60s} scale-down policy, and an identical scale-up
budget; only the signal differed. Operators tempted to tune the HPA
\texttt{behavior} block in an attempt to recover KEDA's observed
responsiveness are therefore optimising a second-order variable.

\paragraph{Why CPU fails on queue-backed workloads.}~For a
CPU-bound consumer governed by a strict \texttt{limits.cpu}, the CPU
signal saturates the moment the queue has any residual work: the
reported utilisation becomes the ratio of limit to request,
independent of the backlog. The HPA can thus neither \emph{estimate}
the deficit while the burst is active nor \emph{release} capacity
while the drain is still in flight. Queue length, in contrast, is a
direct measurement of outstanding work and carries both pieces of
information simultaneously.

\paragraph{When HPA remains reasonable.}~For synchronous HTTP services
without an intermediate queue, CPU remains a plausible scaling signal
and the HPA's simplicity is attractive. For queue consumers, a hybrid
policy is often the pragmatic production choice: let queue length drive
normal elasticity and retain CPU-based limits as a guardrail against
runaway scale-up.

\paragraph{Limitations and next steps.}~The study is confined to a
single-node cluster, one EC2 instance type, one broker, and three
replications per group; it therefore demonstrates a clear directional
effect rather than a universal law. The $\approx\,355$\,s post-burst
window also truncates the Baseline HPA's scale-down tail, so the
observed cost gap is conservative. Future work should sweep
\texttt{cpu} and queue-length targets, compare against HPA with native
\texttt{External} metrics, and repeat the study on multi-node clusters.
